\documentclass[12pt]{article}
\usepackage[utf8]{inputenc}
\usepackage[T5]{fontenc}
\usepackage[vietnamese]{babel}
\usepackage{graphicx,booktabs,hyperref}
\title{Nhận diện trạng thái ô đỗ xe từ ảnh toàn cảnh}
\author{Nhóm nghiên cứu -- Tuấn}
\date{28/10/2026}
\begin{document}
\maketitle
\section{Bài toán và phạm vi}
Dự đoán trạng thái có xe/trống của từng ô. Vùng ô đỗ không phải hộp xe.
\section{Dữ liệu và kiểm tra chất lượng}
PKLot: ảnh toàn cảnh, XML vùng ô, ngày, camera và thời tiết.
Điền số liệu từ audit\_summary.json sau khi tải và kiểm tra toàn bộ dữ liệu.
\section{Phương pháp}
\subsection{Hình học ô đỗ và giao diện detector}
\subsection{Suy luận trạng thái có xe/trống}
\section{Thiết lập đánh giá}
Tách tập theo ngày; occupied là lớp dương. Báo cáo confusion matrix,
accuracy, precision, recall và F1 theo camera và thời tiết.
Không dùng nhãn vùng ô làm ground truth hộp xe để tính mAP.
\section{Kết quả và thảo luận}
Chưa có kết quả mô hình; không dùng dữ liệu kiểm thử nhân tạo làm kết quả thực nghiệm.
\section{Kết luận}
\begin{thebibliography}{9}
\bibitem{pklot} Almeida et al. PKLot -- A robust dataset for parking lot classification.
Expert Systems with Applications, 42(11):4937--4949, 2015.
\end{thebibliography}
\end{document}
